% Options for packages loaded elsewhere
\PassOptionsToPackage{unicode}{hyperref}
\PassOptionsToPackage{hyphens}{url}
\documentclass[
  a4paper,
]{article}
\usepackage{xcolor}
\usepackage[margin=25mm]{geometry}
\usepackage{amsmath,amssymb}
\usepackage{mathrsfs}
\setcounter{secnumdepth}{-\maxdimen} % remove section numbering
\usepackage{iftex}
\ifPDFTeX
  \usepackage[T1]{fontenc}
  \usepackage[utf8]{inputenc}
  \usepackage{textcomp} % provide euro and other symbols
\else % if luatex or xetex
  \usepackage{unicode-math} % this also loads fontspec
  \defaultfontfeatures{Scale=MatchLowercase}
  \defaultfontfeatures[\rmfamily]{Ligatures=TeX,Scale=1}
\fi
\usepackage{lmodern}
\ifPDFTeX\else
  % xetex/luatex font selection
\fi
% Use upquote if available, for straight quotes in verbatim environments
\IfFileExists{upquote.sty}{\usepackage{upquote}}{}
\IfFileExists{microtype.sty}{% use microtype if available
  \usepackage[]{microtype}
  \UseMicrotypeSet[protrusion]{basicmath} % disable protrusion for tt fonts
}{}
\makeatletter
\@ifundefined{KOMAClassName}{% if non-KOMA class
  \IfFileExists{parskip.sty}{%
    \usepackage{parskip}
  }{% else
    \setlength{\parindent}{0pt}
    \setlength{\parskip}{6pt plus 2pt minus 1pt}}
}{% if KOMA class
  \KOMAoptions{parskip=half}}
\makeatother
\usepackage{longtable,booktabs,array}
\newcounter{none} % for unnumbered tables
\usepackage{calc} % for calculating minipage widths
% Correct order of tables after \paragraph or \subparagraph
\usepackage{etoolbox}
\makeatletter
\patchcmd\longtable{\par}{\if@noskipsec\mbox{}\fi\par}{}{}
\makeatother
% Allow footnotes in longtable head/foot
\IfFileExists{footnotehyper.sty}{\usepackage{footnotehyper}}{\usepackage{footnote}}
\makesavenoteenv{longtable}
\setlength{\emergencystretch}{3em} % prevent overfull lines
\providecommand{\tightlist}{%
  \setlength{\itemsep}{0pt}\setlength{\parskip}{0pt}}
\usepackage{bookmark}
\IfFileExists{xurl.sty}{\usepackage{xurl}}{} % add URL line breaks if available
\urlstyle{same}
\hypersetup{
  hidelinks,
  pdfcreator={LaTeX via pandoc}}

\author{}
\date{}

\begin{document}

\section{An Axiomatic Framework for Physical Symmetry Selection by
Anonymous Partial
Projection}\label{an-axiomatic-framework-for-physical-symmetry-selection-by-anonymous-partial-projection}

\subsection{\texorpdfstring{Conditional Identification of \(E_8\) in an
Eight-Component Two-Quartet
Branch}{Conditional Identification of E\_8 in an Eight-Component Two-Quartet Branch}}\label{conditional-identification-of-e_8-in-an-eight-component-two-quartet-branch}

\textbf{Author:} Noriaki Kihara\\
\textbf{ORCID:} 0009-0004-6753-4020\\
\textbf{Date:} July 24, 2026\\
\textbf{Version DOI:}
\href{https://doi.org/10.5281/zenodo.21521900}{10.5281/zenodo.21521900}\\
\textbf{Concept DOI:}
\href{https://doi.org/10.5281/zenodo.21521899}{10.5281/zenodo.21521899}\\
\textbf{Status:} Independent paper, English version v1 / Mathematical
physics and foundations of physics

\begin{center}\rule{0.5\linewidth}{0.5pt}\end{center}

\subsection{Abstract}\label{abstract}

When the solution set of a foundational axiom system contains more than
one closed system, the symmetry observed in physics need not be
identified with the symmetry of the entire solution set. This paper
proposes an \textbf{anonymous partial-projection selection principle},
in which physical state space is defined as a selected
partial-projection image of the candidate space allowed by the
foundational axioms.

Let \(\mathscr S\) be the candidate state space, and let
\(\mathfrak P(\mathscr S)\) be the set of admissible partial projections
satisfying anonymity, scale anonymity, closure consistency, and readout
consistency. The central axiom of this paper is

\[
\boxed{
\exists\mathcal P_*\in\mathfrak P(\mathscr S)
\quad\text{such that}\quad
\mathscr S_{\mathrm{phys}}
{}=
\operatorname{Im}\mathcal P_*.
}
\]

This axiom does not specify which symmetry group is to be selected.
Physical symmetry is identified only after projection, as the
automorphism group of the relational structure carried by the selected
image.

As a nontrivial application, consider a branch in which transition
differences in the selected image lie in an eight-dimensional
positive-definite readout space, split into two four-component systems,
with each local transition lattice of type \(D_4\). Identify the two
\(D_4\) discriminant groups with a single common center and retain only
the neutral sector under its diagonal action. The permitted glue classes
are then restricted to

\[
(0,0),\qquad(v,v),\qquad(s,s),\qquad(c,c).
\]

The lattice obtained by this diagonal gluing is positive definite,
eight-dimensional, even, and unimodular; it is therefore isomorphic to
the \(E_8\) lattice. Its minimal vectors number

\[
48+3\times64=240,
\]

and the dimension of the associated complex simple Lie algebra is

\[
8+240=248.
\]

Furthermore, a Coxeter element of \(E_8\) has order 30.

The paper does not claim that the foundational axioms uniquely force
\(E_8\). The physical reasons for selecting eight components, the
two-quartet \(D_4\) structure, and the common-center projection remain
selection problems. What is established is the axiomatic order:
anonymous partial projection is formulated first, and \(E_8\) appears
only as the mathematical classification of one explicitly conditioned
branch.

\begin{center}\rule{0.5\linewidth}{0.5pt}\end{center}

\subsection{0. Position and Scope}\label{position-and-scope}

\subsubsection{0.1 What this paper
addresses}\label{what-this-paper-addresses}

This paper does three things.

\begin{enumerate}
\def\labelenumi{\arabic{enumi}.}
\tightlist
\item
  It separates the admissible solution space of the foundational axioms
  from the physically realized state space.
\item
  It formulates physical realization as the image of an anonymous
  admissible partial projection.
\item
  It proves that, when a common-center neutrality projection is imposed
  on an eight-component two-quartet branch, the transition lattice of
  the selected image is \(E_8\).
\end{enumerate}

\subsubsection{0.2 What this paper does not
claim}\label{what-this-paper-does-not-claim}

This paper does not claim:

\begin{enumerate}
\def\labelenumi{\arabic{enumi}.}
\tightlist
\item
  that the foundational axioms alone uniquely derive \(E_8\);
\item
  that every physical partial-projection image has \(E_8\) symmetry;
\item
  that a dynamics or selection functional determining the actual
  projection \(\mathcal P_*\) has already been obtained;
\item
  that the eight readout quantities have already been mapped to
  particles, interactions, or quantities of the Standard Model;
\item
  that a finite-recurrence operator is physically identical to a Coxeter
  element of \(E_8\); or
\item
  that \(E_8\) is the symmetry of spacetime as a whole or of the entire
  state space.
\end{enumerate}

\subsubsection{0.3 Central claim}\label{central-claim}

The central claim is an ordering principle:

\[
\boxed{
\begin{aligned}
&\text{foundational axioms}\\
&\Longrightarrow
\text{candidate space containing multiple admissible closed systems}\\
&\Longrightarrow
\text{set of anonymous admissible partial projections}\\
&\Longrightarrow
\text{selection of one physical partial image}\\
&\Longrightarrow
\text{mathematical identification of the symmetry of that image}.
\end{aligned}
}
\]

\(E_8\) is not inserted as an axiom at the beginning of this chain. It
is an identification obtained at the end of one particular branch.

\begin{center}\rule{0.5\linewidth}{0.5pt}\end{center}

\subsection{1. Connection to the Foundational
Axioms}\label{connection-to-the-foundational-axioms}

\subsubsection{1.1 Inherited formation
rules}\label{inherited-formation-rules}

This paper inherits the following formation rules from the \emph{Basic
Axiom System for an Anonymous Equal-Amplitude Composite-Wave Model, v8}
{[}1{]}.

\textbf{Anonymity.} No elementary component is assigned an individual
name, privileged axis, privileged sign, or privileged type. A physical
name or the name of an existing theory may not be used as a reason to
alter an equation or readout rule at the first-principles level.

\textbf{Scale anonymity.} The states \(X\) and \(\lambda X\), for
\(\lambda\in\mathbb C^\ast\), are equivalent as relational structures
after removal of absolute scale.

\textbf{All-positive-sign zero closure:}

\[
Q(X):=\sum_{j=1}^{N}x_j^2=0.
\]

\textbf{Nontrivial existence:}

\[
X\neq0.
\]

\subsubsection{1.2 What is added here}\label{what-is-added-here}

The inherited formation rules restrict candidate states, but they do not
uniquely determine which candidates are physically readable. This paper
fills that logical gap by formulating the \textbf{existence and
admissibility of a partial projection}.

No particular symmetry group is added as an axiom. In particular,

\[
\text{“physical symmetry is }E_8\text{”}
\]

is not postulated. Such a postulate would violate anonymity by using the
desired name of the resulting theory to select the readout rule.

\begin{center}\rule{0.5\linewidth}{0.5pt}\end{center}

\subsection{2. Admissible Solution
Space}\label{admissible-solution-space}

\subsubsection{2.1 Projectivized candidate
states}\label{projectivized-candidate-states}

Define scale equivalence by

\[
X\sim_{\mathrm{sc}}\lambda X,
\qquad
\lambda\in\mathbb C^\ast.
\]

The projectivized set of nontrivial zero-closure candidates is

\[
\mathscr S_0
:=
\left\{
[X]
\ \middle|\
X\neq0,\ Q(X)=0
\right\},
\]

where

\[
[X]
:=
\left\{
\lambda X
\mid
\lambda\in\mathbb C^\ast
\right\}.
\]

\subsubsection{2.2 Candidate systems with
evolution}\label{candidate-systems-with-evolution}

Let \(U\) be a complex-linear evolution operator on the state space.
Linearity makes \(U\) well defined on scale-equivalence classes.
Compatibility with the same closed system is expressed by

\[
Q(UY)=Q(Y)
\qquad
(\forall Y).
\]

Define the candidate systems with evolution by

\[
\mathscr S
:=
\left\{
([X],U)
\ \middle|\
[X]\in\mathscr S_0,\ Q\circ U=Q
\right\}.
\]

\subsubsection{2.3 Finite-recurrence
branch}\label{finite-recurrence-branch}

Candidate systems with finite recurrence form

\[
\mathscr S_{\mathrm{fin}}
:=
\left\{
([X],U,n)
\ \middle|\
([X],U)\in\mathscr S,\quad
U^n=I,\quad
n\in\mathbb N
\right\}.
\]

Finite recurrence is not required for the general partial-projection
principle. It is an additional condition selecting the subfamily that
may later be compared with a Coxeter period of \(E_8\).

\begin{center}\rule{0.5\linewidth}{0.5pt}\end{center}

\subsection{3. Anonymous Partial
Projection}\label{anonymous-partial-projection}

\subsubsection{3.1 Definitions}\label{definitions}

\textbf{Definition 3.1 (partial projection).} A partial projection of
the candidate space \(\mathscr S\) is a pair

\[
\mathcal P
=
\left(
\mathscr D_{\mathcal P},
p_{\mathcal P}
\right),
\]

where \(\mathscr D_{\mathcal P}\subseteq\mathscr S\) and

\[
p_{\mathcal P}:
\mathscr D_{\mathcal P}
\longrightarrow
\mathscr Y_{\mathcal P}.
\]

Here ``partial projection'' is not restricted to an orthogonal
projection in linear algebra. It may include:

\begin{enumerate}
\def\labelenumi{\arabic{enumi}.}
\tightlist
\item
  selecting only part of the candidate systems as its domain;
\item
  quotienting differences that are unreadable; and
\item
  mapping readable relational quantities into another representation
  space.
\end{enumerate}

\textbf{Definition 3.2 (nontrivial partial projection).} A partial
projection is nontrivial if at least one of the following holds:

\[
\mathscr D_{\mathcal P}
\subsetneq
\mathscr S,
\]

or

\[
\exists s_1\neq s_2
\quad\text{such that}\quad
p_{\mathcal P}(s_1)
=
p_{\mathcal P}(s_2).
\]

Thus either some candidates are removed, or some differences between
candidates are identified as unreadable. The identity map alone is not
counted as a nontrivial realization of the principle.

\subsubsection{3.2 Transition-difference
lattice}\label{transition-difference-lattice}

Assume that the selected image contains discrete transition differences.
Let \(\Delta_{\mathcal P}\subseteq V_{\mathcal P}\) be their set, and
define its integer span by

\[
\Lambda_{\mathcal P}
:=
\operatorname{span}_{\mathbb Z}
\Delta_{\mathcal P}.
\]

\(\Lambda_{\mathcal P}\) stores readable state differences or transition
differences, not individual names of states.

\subsubsection{3.3 Admissibility
conditions}\label{admissibility-conditions}

\textbf{Definition 3.3 (admissible partial projection).} A partial
projection is admissible if it is nontrivial and satisfies the following
conditions.

\paragraph{A. Component anonymity}\label{a.-component-anonymity}

For any permutation or anonymity-preserving transformation \(g\), the
projection result is independent of the relabeling of individual
components:

\[
p_{\mathcal P}(g\cdot[X])
\sim_{\mathscr Y}
p_{\mathcal P}([X]).
\]

\paragraph{B. Scale anonymity}\label{b.-scale-anonymity}

\[
p_{\mathcal P}([\lambda X])
=
p_{\mathcal P}([X]),
\qquad
\lambda\in\mathbb C^\ast.
\]

\paragraph{C. Closure consistency}\label{c.-closure-consistency}

The projection domain consists of candidates satisfying the foundational
closure, and the projected transition differences are closed under the
composition rule defined on the image.

\paragraph{D. Readout consistency}\label{d.-readout-consistency}

Two states identified by the projection cannot be distinguished using
only the adopted readout quantities.

\paragraph{E. Prohibition of prior physical
naming}\label{e.-prohibition-of-prior-physical-naming}

Neither a physical name nor a symmetry-group name that becomes available
only after projection may be used in the definition of the projection.

Write the set of all admissible partial projections as

\[
\mathfrak P(\mathscr S).
\]

\begin{center}\rule{0.5\linewidth}{0.5pt}\end{center}

\subsection{4. Partial-Projection Existence
Axiom}\label{partial-projection-existence-axiom}

\subsubsection{4.1 The axiom}\label{the-axiom}

\textbf{Axiom PP (anonymous partial-projection existence axiom).} For
the candidate space \(\mathscr S\) of the foundational axioms, the
physically realized state space is the image of at least one nontrivial
admissible partial projection:

\[
\boxed{
\exists\mathcal P_*
\in
\mathfrak P(\mathscr S)
\quad\text{such that}\quad
\mathscr S_{\mathrm{phys}}
=
\operatorname{Im}\mathcal P_*.
}
\]

\subsubsection{4.2 What the axiom does not
assert}\label{what-the-axiom-does-not-assert}

Axiom PP does not assert that

\[
\mathcal P_*
\text{ is unique}.
\]

Nor does it specify a selection functional such as

\[
\mathcal P_*
=
\operatorname*{arg\,max}_{\mathcal P\in\mathfrak P(\mathscr S)}
\mathcal F(\mathcal P).
\]

It is therefore an existence axiom saying that physical realization is a
partial-projection image. It is not yet a selection dynamics determining
which image is realized.

\subsubsection{4.3 Physical symmetry}\label{physical-symmetry}

\textbf{Definition 4.1 (physical symmetry of a projection image).}
Define the physical symmetry group to be the automorphism group
preserving the readout relations, transition-difference set, and
composition rule of the selected image:

\[
G_{\mathrm{phys}}
:=
\operatorname{Aut}
\left(
\operatorname{Im}\mathcal P_*,
\Delta_{\mathcal P_*}
\right).
\]

For a description by a discrete transition lattice, write

\[
G_{\mathrm{lattice}}
:=
\operatorname{Aut}
\left(
\Lambda_{\mathcal P_*},
(\,\cdot\,,\,\cdot\,)
\right),
\]

where a lattice automorphism must preserve both the lattice and its
inner product.

Physical symmetry is not a name assigned in advance to the full
candidate space. It is identified after projection as the group
preserving the relations that survive in the selected image.

\subsubsection{4.4 Nonunique branches}\label{nonunique-branches}

Axiom PP does not exclude two admissible projections

\[
\mathcal P_1,\mathcal P_2
\in
\mathfrak P(\mathscr S)
\]

for which

\[
\operatorname{Im}\mathcal P_1
\not\cong
\operatorname{Im}\mathcal P_2.
\]

Different stability conditions, readout conditions, or boundary
conditions may therefore realize different symmetry branches.

\begin{center}\rule{0.5\linewidth}{0.5pt}\end{center}

\subsection{5. Eight-Component Two-Quartet
Branch}\label{eight-component-two-quartet-branch}

We now impose additional conditions on one branch allowed by Axiom PP.
They are not consequences of PP; they are explicit branch assumptions
defining the \(E_8\) branch.

\subsubsection{5.1 Branch assumption Q1: positive-definite
eight-component
readout}\label{branch-assumption-q1-positive-definite-eight-component-readout}

Assume that the transition differences of the selected image lie in an
eight-dimensional real readout space with a positive-definite inner
product,

\[
V_{\mathrm{read}}
\cong
\mathbb R^8,
\]

and that

\[
\operatorname{rank}
\Lambda_{\mathcal P_*}
=8.
\]

The complex bilinear form inherited from Axiom 1,

\[
Q(X)=\sum_jx_j^2,
\]

and the positive-definite readout inner product

\[
(\alpha,\beta)
\]

introduced here are different structures. Positive definiteness is not
derived directly from \(Q(X)=0\).

\subsubsection{5.2 Branch assumption Q2: two-quartet
decomposition}\label{branch-assumption-q2-two-quartet-decomposition}

Assume an orthogonal decomposition of the readout space into two
four-component subspaces:

\[
V_{\mathrm{read}}
=
V_A\oplus V_B,
\qquad
\dim V_A
=
\dim V_B
=4.
\]

At this stage no component is named space, time, mass, charge, or any
other physical quantity.

\subsubsection{\texorpdfstring{5.3 Branch assumption Q3: local \(D_4\)
transition
lattices}{5.3 Branch assumption Q3: local D\_4 transition lattices}}\label{branch-assumption-q3-local-d_4-transition-lattices}

Assume that the local transition lattice of each four-component
subsystem is isomorphic to the lattice defined by

\[
D_4
:=
\left\{
z\in\mathbb Z^4
\ \middle|\
\sum_{i=1}^{4}z_i
\equiv0
\pmod2
\right\}.
\]

The symbol \(D_4\) is used as its mathematical classification name. The
content of the branch assumption is the displayed coordinate condition;
the lattice is not selected because a symmetry name was desired.

Before gluing, the local lattice is therefore

\[
L_0
:=
D_4^{(A)}
\oplus
D_4^{(B)}.
\]

This local \(D_4\) structure has not been derived from the foundational
axioms or Axiom PP alone.

\subsubsection{\texorpdfstring{5.4 Discriminant classes of
\(D_4\)}{5.4 Discriminant classes of D\_4}}\label{discriminant-classes-of-d_4}

Let \(D_4^*\) be the dual lattice of \(D_4\). Its discriminant group is

\[
\mathcal A
:=
D_4^*/D_4
\cong
(\mathbb Z/2\mathbb Z)^2,
\]

with four classes

\[
\mathcal A
=
\{0,v,s,c\}.
\]

Using a standard orthonormal basis \(e_1,\ldots,e_4\), representatives
may be chosen as

\[
v=e_1+D_4,
\]

\[
s=
\frac12
(e_1+e_2+e_3+e_4)
+D_4,
\]

and

\[
c=
\frac12
(-e_1+e_2+e_3+e_4)
+D_4.
\]

The discriminant quadratic form

\[
q:
\mathcal A
\longrightarrow
\mathbb Q/2\mathbb Z,
\qquad
q(a)
=
(a,a)
\bmod
2\mathbb Z
\]

satisfies

\[
q(v)
=
q(s)
=
q(c)
=1
\pmod{2\mathbb Z}.
\]

\begin{center}\rule{0.5\linewidth}{0.5pt}\end{center}

\subsection{6. Common-Center Neutrality
Projection}\label{common-center-neutrality-projection}

\subsubsection{6.1 Branch assumption Q4: common-center
action}\label{branch-assumption-q4-common-center-action}

Assume that the discriminant groups of the two four-component systems
are not retained as independent label groups. Instead, identify them
with one common center label group

\[
\mathcal A
\cong
(\mathbb Z/2\mathbb Z)^2,
\]

on which the same element \(z\in\mathcal A\) acts. Only the finite
Abelian group structure of \(\mathcal A\) and its discriminant pairing
are used here.

The nondegenerate pairing induced by the discriminant form identifies
\(\mathcal A\) with its character group \(\widehat{\mathcal A}\). Let
\(\chi_a\) denote the central character corresponding to
\(a\in\mathcal A\). Every character takes real values \(\pm1\). Write
the two discriminant classes as

\[
(a,b)
\in
\mathcal A\oplus\mathcal A.
\]

\subsubsection{6.2 Branch assumption Q5: selection of the neutral
sector}\label{branch-assumption-q5-selection-of-the-neutral-sector}

Assume that only components invariant under the common-center action
remain in the physical readout image. Define the neutrality projector by

\[
\Pi_{\mathrm{com}}(a,b)
:=
\frac{1}{|\mathcal A|}
\sum_{z\in\mathcal A}
\chi_a(z)\chi_b(z).
\]

The orthogonality relation for characters of a finite Abelian group
gives

\[
\Pi_{\mathrm{com}}(a,b)
=
\delta_{a+b,\,0}.
\]

\subsubsection{6.3 Diagonal-class selection
theorem}\label{diagonal-class-selection-theorem}

\textbf{Theorem 6.1 (diagonal-class selection by a common center).}
Under branch assumptions Q4 and Q5, the only discriminant classes
passing the common-center neutrality projection are

\[
\boxed{
(0,0),
\qquad
(v,v),
\qquad
(s,s),
\qquad
(c,c).
}
\]

\textbf{Proof.} Neutrality requires

\[
a+b=0.
\]

Every element of \(\mathcal A\cong(\mathbb Z/2\mathbb Z)^2\) has order
one or two, so

\[
-a=a.
\]

Hence

\[
b=-a=a.
\]

Substitution of \(\mathcal A=\{0,v,s,c\}\) leaves precisely the four
displayed classes. \(\square\)

\subsubsection{6.4 Diagonal glue group}\label{diagonal-glue-group}

Define the admissible glue group by

\[
H_{\mathrm{diag}}
:=
\left\{
(0,0),
(v,v),
(s,s),
(c,c)
\right\}.
\]

A graph gluing obtained by simultaneously relabeling \(v,s,c\) on one
side by triality can be written in the same diagonal form after the
identification of the two discriminant groups is relabeled. Thus ``equal
classes'' is the standard form after both discriminant groups have been
identified with the same physical center labels.

\begin{center}\rule{0.5\linewidth}{0.5pt}\end{center}

\subsection{\texorpdfstring{7. Conditional Identification of the \(E_8\)
Lattice}{7. Conditional Identification of the E\_8 Lattice}}\label{conditional-identification-of-the-e_8-lattice}

\subsubsection{7.1 Glued lattice}\label{glued-lattice}

Define the lattice after diagonal gluing by

\[
\Lambda
:=
\bigcup_{h\in H_{\mathrm{diag}}}
\left(
L_0+h
\right).
\]

The general theory of lattice gluing and discriminant forms follows
Nikulin {[}2{]}. The index, determinant, and parity needed here are also
calculated directly below.

\subsubsection{7.2 Index and determinant}\label{index-and-determinant}

Since the determinant of \(D_4\) is 4,

\[
\det L_0
=
\det(D_4\oplus D_4)
=
4^2
=16.
\]

The group \(H_{\mathrm{diag}}\) has order 4, so

\[
[\Lambda:L_0]
=4.
\]

The determinant formula for a finite-index overlattice gives

\[
\det\Lambda
=
\frac{\det L_0}
{[\Lambda:L_0]^2}
=
\frac{16}{4^2}
=1.
\]

Thus \(\Lambda\) is unimodular.

\subsubsection{7.3 Evenness}\label{evenness}

Let \(q\oplus q\) be the direct-sum discriminant form. For every nonzero
diagonal class,

\[
(q\oplus q)(a,a)
=
q(a)+q(a)
=
1+1
=0
\pmod{2\mathbb Z},
\]

where

\[
a\in\{v,s,c\}.
\]

Therefore \(H_{\mathrm{diag}}\) is isotropic with respect to the
discriminant form, and the glued lattice \(\Lambda\) is even.

\subsubsection{7.4 Main theorem}\label{main-theorem}

\textbf{Theorem 7.1 (\(E_8\) identification in the eight-component
two-quartet branch).} Under branch assumptions Q1--Q5, the
transition-difference lattice \(\Lambda\) obtained by the common-center
neutrality projection is isomorphic to the \(E_8\) lattice:

\[
\boxed{
\Lambda
\cong
\Lambda_{E_8}.
}
\]

\textbf{Proof.} By Q1 and Q2, \(\Lambda\) is positive definite and
eight-dimensional. By Q3, the lattice before gluing is
\(D_4\oplus D_4\). By Theorem 6.1, the glue group is
\(H_{\mathrm{diag}}\). Section 7.2 gives \(\det\Lambda=1\), and Section
7.3 shows that \(\Lambda\) is even. Thus \(\Lambda\) is a
positive-definite, rank-eight, even unimodular lattice. In dimension
eight, this lattice is unique up to isomorphism and is the \(E_8\)
lattice {[}3{]}. Hence \(\Lambda\cong\Lambda_{E_8}\). \(\square\)

\subsubsection{7.5 Logical status}\label{logical-status}

Theorem 7.1 is not the unconditional statement

\[
Q(X)=0
\Longrightarrow
E_8.
\]

Its exact content is

\[
\boxed{
\begin{array}{c}
\text{Axiom PP}\\
{}+\text{Q1: positive-definite rank eight}\\
{}+\text{Q2: two quartets}\\
{}+\text{Q3: local }D_4\\
{}+\text{Q4--Q5: common-center neutrality projection}
\end{array}
\Longrightarrow
\Lambda_{E_8}.
}
\]

\(E_8\) is therefore not a name inserted into the axioms. It is the
classification of a projection image satisfying explicit conditions.

\begin{center}\rule{0.5\linewidth}{0.5pt}\end{center}

\subsection{8. The 240 Roots and Dimension
248}\label{the-240-roots-and-dimension-248}

\subsubsection{8.1 Minimal transition
vectors}\label{minimal-transition-vectors}

Let the norm-two vectors of the \(E_8\) lattice be

\[
\Phi
:=
\left\{
\alpha\in\Lambda
\mid
(\alpha,\alpha)=2
\right\}.
\]

\subsubsection{\texorpdfstring{8.2 Vectors inside
\(D_4\oplus D_4\)}{8.2 Vectors inside D\_4\textbackslash oplus D\_4}}\label{vectors-inside-d_4oplus-d_4}

The roots of \(D_4\) are

\[
\pm e_i\pm e_j,
\qquad
1\le i<j\le4,
\]

and their number is

\[
4\binom42
=24.
\]

Hence the number of norm-two vectors inside

\[
D_4\oplus D_4
\]

is

\[
24+24=48.
\]

\subsubsection{8.3 Nonzero glue classes}\label{nonzero-glue-classes}

Each of the discriminant classes \(v,s,c\) has eight minimal
representatives of norm 1. Each diagonal class

\[
(v,v),\qquad(s,s),\qquad(c,c)
\]

therefore contributes

\[
8\times8=64
\]

norm-two vectors. Since there are three nonzero diagonal classes,

\[
\boxed{
|\Phi|
=
48+3\times64
=240.
}
\]

\subsubsection{8.4 Closure under
reflections}\label{closure-under-reflections}

For each \(\alpha\in\Phi\), define

\[
s_\alpha(x)
:=
x
-\frac{2(x,\alpha)}
{(\alpha,\alpha)}
\alpha
=
x-(x,\alpha)\alpha.
\]

Because \(\Lambda\) is integral and \((\alpha,\alpha)=2\),

\[
s_\alpha(\Lambda)
=
\Lambda.
\]

The group generated by these reflections is the Weyl group \(W(E_8)\) of
the \(E_8\) root system.

\subsubsection{8.5 Lie algebra}\label{lie-algebra}

Using an eight-dimensional Cartan subalgebra \(\mathfrak h\) and a
one-dimensional root space \(\mathfrak g_\alpha\) for every root,

\[
\mathfrak e_8(\mathbb C)
=
\mathfrak h
\oplus
\bigoplus_{\alpha\in\Phi}
\mathfrak g_\alpha.
\]

Consequently,

\[
\boxed{
\dim\mathfrak e_8
=
8+240
=248.
}
\]

The reflection symmetry of the discrete transition lattice is
\(W(E_8)\), while the complex simple Lie algebra containing the root
spaces is \(\mathfrak e_8(\mathbb C)\). If the compact real form
corresponding to the positive-definite inner product is chosen, the
associated simply connected Lie group is compact \(E_8\).

\begin{center}\rule{0.5\linewidth}{0.5pt}\end{center}

\subsection{9. Connection to Finite
Recurrence}\label{connection-to-finite-recurrence}

\subsubsection{9.1 Coxeter element}\label{coxeter-element}

The Coxeter number of the \(E_8\) root system is 30, and a Coxeter
element \(C\) satisfies {[}4{]}

\[
\boxed{
C^{30}=I.
}
\]

Its eigenvalues on the Cartan space are

\[
\operatorname{Spec}(C)
=
\left\{
\exp
\left(
\frac{2\pi i m}{30}
\right)
\ \middle|\
m\in
\{1,7,11,13,17,19,23,29\}
\right\}.
\]

Under the Coxeter action, the 240 roots split into eight orbits of
length 30 {[}4{]}:

\[
\Phi
=
\bigsqcup_{a=1}^{8}
\mathcal O_a,
\qquad
|\mathcal O_a|
=30.
\]

Thus

\[
\boxed{
248
=
8+240
=
8+8\times30
=
8(30+1).
}
\]

\subsubsection{9.2 Logical relation to a finite-recurrence
axiom}\label{logical-relation-to-a-finite-recurrence-axiom}

An \(E_8\) projection image contains an order-30 finite-recurrence
operator \(C\). It does not follow from Theorem 7.1, however, that a
general physical evolution operator \(U\) satisfying

\[
U^n=I
\]

is identical to \(C\):

\[
U=C.
\]

Identifying

\[
U_{\mathrm{phys}}
\longleftrightarrow
C
\]

requires an additional correspondence showing that physical evolution
cycles through the complete root system along Coxeter orbits.

The result established here is only

\[
\boxed{
\Lambda_{\mathcal P_*}
\cong
\Lambda_{E_8}
\Longrightarrow
\exists C\in\operatorname{Aut}(\Lambda_{\mathcal P_*})
\text{ such that }
C^{30}=I.
}
\]

\begin{center}\rule{0.5\linewidth}{0.5pt}\end{center}

\subsection{10. Consistency with
Anonymity}\label{consistency-with-anonymity}

\subsubsection{10.1 No symmetry-group name in the projection
condition}\label{no-symmetry-group-name-in-the-projection-condition}

Axiom PP and the admissibility conditions contain no reference to the
name \(E_8\). The branch assumptions are stated as:

\begin{enumerate}
\def\labelenumi{\arabic{enumi}.}
\tightlist
\item
  eight components;
\item
  two four-component systems;
\item
  an even-integer condition on local transition differences; and
\item
  a common-center neutrality projection.
\end{enumerate}

The resulting lattice is then identified as \(E_8\) by the known
classification of lattices. The logical order is therefore

\[
\boxed{
\text{projection conditions}
\longrightarrow
\text{lattice structure}
\longrightarrow
\text{symmetry-group name}.
}
\]

\subsubsection{10.2 Physical names are also assigned after
projection}\label{physical-names-are-also-assigned-after-projection}

A later interpretation of the eight components as, for example,

\[
(x,y,z,R)
\oplus
(t,Q_1,Q_2,Q_3)
\]

is not part of the theorem. Such a correspondence is a physical
interpretation that can be made only after readout operations, conserved
quantities, exchange rules, and experimental correspondences have been
specified for the components.

\subsubsection{\texorpdfstring{10.3 Branches other than
\(E_8\)}{10.3 Branches other than E\_8}}\label{branches-other-than-e_8}

If any of Q1--Q5 fails, this paper does not require \(E_8\).

\begin{itemize}
\tightlist
\item
  If the readout rank is not eight, Theorem 7.1 does not apply.
\item
  If the local lattice is not of type \(D_4\), a different gluing
  problem results.
\item
  Without diagonal gluing, \(D_4\oplus D_4\) remains a decomposable
  lattice of determinant 16.
\item
  Without positive definiteness, uniqueness of the positive-definite
  \(E_8\) lattice does not apply.
\end{itemize}

The framework therefore permits other partial-projection images and
other symmetry branches from the outset.

\begin{center}\rule{0.5\linewidth}{0.5pt}\end{center}

\subsection{11. The Selection Problem}\label{the-selection-problem}

\subsubsection{11.1 The unresolved core}\label{the-unresolved-core}

The unresolved part of the framework is not the lattice calculation
after gluing. It is why an actual physical system selects a projection
satisfying Q1--Q5.

What remains necessary is

\[
\boxed{
\text{an anonymous selection rule for }\mathcal P_*.
}
\]

\subsubsection{11.2 Representation by a selection
functional}\label{representation-by-a-selection-functional}

If the selection rule can be expressed as a functional, introduce

\[
\mathcal F:
\mathfrak P(\mathscr S)
\longrightarrow
\mathbb R
\]

and write

\[
\mathcal P_*
\in
\operatorname*{arg\,ext}_{\mathcal P\in\mathfrak P(\mathscr S)}
\mathcal F(\mathcal P).
\]

Here \(\operatorname*{arg\,ext}\) denotes a maximum, minimum, or
stationarity condition. No specific form of \(\mathcal F\) is assumed in
this paper.

\subsubsection{11.3 Requirements on the selection
rule}\label{requirements-on-the-selection-rule}

At minimum, the selection rule must:

\begin{enumerate}
\def\labelenumi{\arabic{enumi}.}
\tightlist
\item
  not refer to individual component names;
\item
  not refer to an absolute scale;
\item
  preserve closure;
\item
  not insert a symmetry-group name directly into its objective; and
\item
  make the selection of eight components, local \(D_4\) structure, and
  common-center action comparable under one criterion.
\end{enumerate}

\subsubsection{11.4 Status of the common-center
condition}\label{status-of-the-common-center-condition}

The common-center neutrality projection is a sufficient condition for
diagonal gluing. However,

\[
\text{why the same center acts on both quartets}
\]

and

\[
\text{why only the neutral sector is readable}
\]

must be derived from a selection dynamics or an observational
construction.

The common-center condition is therefore currently a \textbf{candidate
selection axiom} for the \(E_8\) branch, not a derived consequence of
the foundational axioms.

\begin{center}\rule{0.5\linewidth}{0.5pt}\end{center}

\subsection{12. Discriminating Tests}\label{discriminating-tests}

The framework calls for the following calculations.

\subsubsection{12.1 Readout rank}\label{readout-rank}

Test whether the rank of physically independent transition differences
is

\[
\operatorname{rank}\Lambda_{\mathcal P_*}
=8.
\]

\subsubsection{12.2 Local lattice}\label{local-lattice}

Test whether the minimal transition vectors of each four-component
system close under the \(D_4\) root condition

\[
\pm e_i\pm e_j.
\]

\subsubsection{12.3 Correlation of discriminant
classes}\label{correlation-of-discriminant-classes}

Test whether simultaneous occurrence of the discriminant classes in the
two four-component systems satisfies

\[
(a,b)
\text{ occurs}
\quad\Longleftrightarrow\quad
a=b.
\]

If off-diagonal classes such as

\[
(v,s),\qquad(v,c),\qquad(s,c)
\]

remain stable, the common-center neutrality-projection hypothesis is
rejected.

\subsubsection{12.4 Coxeter period}\label{coxeter-period}

Test whether all selected transition differences split into eight orbits
of length 30:

\[
240
\stackrel{?}{=}
8\times30.
\]

Even if this test succeeds, identifying physical time evolution \(U\)
with the Coxeter element \(C\) still requires an operator-level
correspondence.

\begin{center}\rule{0.5\linewidth}{0.5pt}\end{center}

\subsection{13. Form of Incorporation into the Basic Axiom
System}\label{form-of-incorporation-into-the-basic-axiom-system}

The general principle to be incorporated into the basic axiom system is
Axiom PP, not \(E_8\) itself.

\subsubsection{13.1 Proposed minimal
axiom}\label{proposed-minimal-axiom}

\begin{quote}
\textbf{Anonymous partial-projection existence axiom.} For the
admissible solution space of the foundational axioms, the physically
realized state space is the image of at least one nontrivial admissible
partial projection satisfying anonymity, scale anonymity, closure
consistency, and readout consistency.
\end{quote}

\[
\mathscr S_{\mathrm{phys}}
=
\operatorname{Im}\mathcal P_*,
\qquad
\mathcal P_*
\in
\mathfrak P(\mathscr S).
\]

\subsubsection{13.2 Items not to be incorporated
directly}\label{items-not-to-be-incorporated-directly}

The following are not inserted directly into the foundational axioms:

\begin{enumerate}
\def\labelenumi{\arabic{enumi}.}
\tightlist
\item
  the symmetry-group name \(E_8\);
\item
  the derived numbers 240, 248, and 30; and
\item
  a declaration that the eight-component two-quartet branch is the
  unique realization.
\end{enumerate}

They remain theorems and consequences of this paper under additional
branch conditions.

\subsubsection{13.3 Consistency with a single source of
truth}\label{consistency-with-a-single-source-of-truth}

The published basic axiom system v8 {[}1{]} remains unchanged. After DOI
publication of the present paper, only the general Axiom PP is to be
added to the next version based on v8. The order is

\[
\begin{aligned}
&\text{fix definitions and theorems in an independent paper}\\
&\longrightarrow\\
&\text{incorporate only the general axiom into the next basic axiom system}.
\end{aligned}
\]

\begin{center}\rule{0.5\linewidth}{0.5pt}\end{center}

\subsection{14. Conclusion}\label{conclusion}

This paper has presented an axiomatic framework in which physical
symmetry is not assigned in advance to the entire solution space of the
foundational axioms, but is identified after projection as an
automorphism of an admissible partial-projection image.

The central axiom is

\[
\boxed{
\mathscr S_{\mathrm{phys}}
=
\operatorname{Im}\mathcal P_*,
\qquad
\mathcal P_*
\in
\mathfrak P(\mathscr S).
}
\]

It does not specify a symmetry group. It states only that a
partial-projection image consistent with anonymity, scale anonymity,
closure, and readout gives the physical realization.

For one particular branch, positive-definite rank-eight readout, a
two-quartet decomposition, local \(D_4\) transition lattices, and a
common-center neutrality projection were assumed. The only surviving
classes are then

\[
\boxed{
(0,0),
\qquad
(v,v),
\qquad
(s,s),
\qquad
(c,c),
}
\]

and the glued lattice satisfies

\[
\boxed{
\Lambda
\cong
\Lambda_{E_8}.
}
\]

It follows that

\[
\boxed{
|\Phi|
=240,
\qquad
\dim\mathfrak e_8
=248,
\qquad
C^{30}
=I.
}
\]

The logical position of \(E_8\) in this paper is therefore

\[
\boxed{
\begin{aligned}
&\text{$E_8$ was not selected by name;}\\
&\text{it is the classification of the image selected}\\
&\text{by an anonymous partial projection.}
\end{aligned}
}
\]

The central remaining problem is to derive an anonymous selection
functional or stability condition that selects the actual
\(\mathcal P_*\) from the set of admissible partial projections.

\begin{center}\rule{0.5\linewidth}{0.5pt}\end{center}

\subsection{References}\label{references}

{[}1{]} N. Kihara, ``Basic Axiom System for an Anonymous Equal-Amplitude
Composite-Wave Model, v8,'' Zenodo, Version DOI:
\href{https://doi.org/10.5281/zenodo.21495422}{10.5281/zenodo.21495422},
Concept DOI:
\href{https://doi.org/10.5281/zenodo.21315735}{10.5281/zenodo.21315735},
2026.

{[}2{]} V. V. Nikulin, ``Integral symmetric bilinear forms and some of
their applications,'' \emph{Mathematics of the USSR-Izvestiya}, vol.~14,
no. 1, pp.~103--167, 1980. DOI:
\href{https://doi.org/10.1070/IM1980v014n01ABEH001060}{10.1070/IM1980v014n01ABEH001060}.

{[}3{]} L. J. Mordell, ``The definite quadratic forms in eight variables
with determinant unity,'' \emph{Journal de Mathématiques Pures et
Appliquées}, 9e série, vol.~17, pp.~41--46, 1938.
\href{https://www.numdam.org/item/JMPA_1938_9_17_1-4_41_0/}{NUMDAM}.

{[}4{]} D. A. Richter, ``Triacontagonal coordinates for the \(E_8\) root
system,'' arXiv:0704.3091, 2007.
\href{https://arxiv.org/abs/0704.3091}{arXiv}.

\begin{center}\rule{0.5\linewidth}{0.5pt}\end{center}

\subsection{Appendix A: Audit of Claim
Status}\label{appendix-a-audit-of-claim-status}

{\def\LTcaptype{none} % do not increment counter
\begin{longtable}[]{@{}
  >{\raggedright\arraybackslash}p{(\linewidth - 4\tabcolsep) * \real{0.3333}}
  >{\raggedright\arraybackslash}p{(\linewidth - 4\tabcolsep) * \real{0.3333}}
  >{\raggedright\arraybackslash}p{(\linewidth - 4\tabcolsep) * \real{0.3333}}@{}}
\toprule\noalign{}
\begin{minipage}[b]{\linewidth}\raggedright
Item
\end{minipage} & \begin{minipage}[b]{\linewidth}\raggedright
Status
\end{minipage} & \begin{minipage}[b]{\linewidth}\raggedright
Treatment in this paper
\end{minipage} \\
\midrule\noalign{}
\endhead
\bottomrule\noalign{}
\endlastfoot
Anonymity, scale anonymity, and zero closure & Existing axioms &
Inherited from the basic axiom system v8 {[}1{]} \\
Physical realization is an admissible partial-projection image & New
Axiom PP & Central axiom of this paper \\
Multiple admissible partial projections may exist & Structure allowed by
PP & Uniqueness is not assumed \\
Rule selecting the actual \(\mathcal P_*\) & Not derived & A selection
functional or dynamics is required \\
Positive-definite rank-eight readout & Branch assumption Q1 & Not
derived from the foundational axioms \\
Two-quartet decomposition & Branch assumption Q2 & Not derived from the
foundational axioms \\
Local \(D_4\) transition lattice & Branch assumption Q3 & Not derived
from the foundational axioms \\
Common-center action & Branch assumption Q4 & Candidate selection
axiom \\
Retention of the neutral sector only & Branch assumption Q5 & Candidate
selection axiom \\
Only diagonal classes survive & Theorem 6.1 & Derived from Q4--Q5 \\
Even unimodular rank-eight lattice & Derived consequence & Derived from
the glue index and discriminant form \\
\(\Lambda\cong\Lambda_{E_8}\) & Conditional Theorem 7.1 & Holds under
Q1--Q5 \\
240 roots & Derived consequence & Direct count by glue classes \\
Dimension 248 & Derived consequence & Rank 8 plus 240 roots \\
\(C^{30}=I\) & Known mathematical consequence after \(E_8\)
identification & Richter {[}4{]} \\
Identification of physical \(U\) with Coxeter element \(C\) & Connection
problem & Requires an operator correspondence \\
Mapping of eight components to \(xyzR,tQ_1Q_2Q_3\) & Physical
interpretation, not derived & Not part of the theorem \\
\end{longtable}
}

\begin{center}\rule{0.5\linewidth}{0.5pt}\end{center}

\subsection{Appendix B: Branch Audit}\label{appendix-b-branch-audit}

{\def\LTcaptype{none} % do not increment counter
\begin{longtable}[]{@{}
  >{\raggedright\arraybackslash}p{(\linewidth - 2\tabcolsep) * \real{0.5000}}
  >{\raggedright\arraybackslash}p{(\linewidth - 2\tabcolsep) * \real{0.5000}}@{}}
\toprule\noalign{}
\begin{minipage}[b]{\linewidth}\raggedright
Condition removed or changed
\end{minipage} & \begin{minipage}[b]{\linewidth}\raggedright
Result
\end{minipage} \\
\midrule\noalign{}
\endhead
\bottomrule\noalign{}
\endlastfoot
Remove Axiom PP & No principle distinguishes candidate space from
physically realized space \\
Remove Q1 & The positive-definite rank-eight classification theorem
cannot be applied \\
Remove Q2 & The construction cannot be expressed as a \(D_4\oplus D_4\)
gluing \\
Remove Q3 & The discriminant classes \(\{0,v,s,c\}\) and the starting
point of diagonal gluing are lost \\
Remove Q4 & The two quartet discriminant classes may be selected
independently \\
Remove Q5 & There is no projection principle excluding off-diagonal
classes \\
Remove diagonal gluing & \(D_4\oplus D_4\) retains determinant 16 and is
not self-dual \\
Replace positive definiteness by indefinite signature & Uniqueness of
the positive-definite \(E_8\) lattice no longer applies \\
Do not assume physical \(U=C\) & An order-30 action exists inside
\(E_8\), but its identity with physical time evolution remains
undetermined \\
\end{longtable}
}

\begin{center}\rule{0.5\linewidth}{0.5pt}\end{center}

\subsection{Appendix C: Minimal Derivation
Chain}\label{appendix-c-minimal-derivation-chain}

\[
\boxed{
\begin{aligned}
&\text{anonymity, scale anonymity, and zero closure}\\
&\Longrightarrow
\mathscr S\\
&\xrightarrow{\text{Axiom PP}}
\operatorname{Im}\mathcal P_*\\
&\xrightarrow{\text{Q1--Q3}}
D_4\oplus D_4\\
&\xrightarrow{\text{Q4--Q5}}
H_{\mathrm{diag}}
=
\{(0,0),(v,v),(s,s),(c,c)\}\\
&\Longrightarrow
\text{positive-definite rank-eight even unimodular lattice}\\
&\Longrightarrow
\Lambda_{E_8}\\
&\Longrightarrow
240\text{ roots}\\
&\Longrightarrow
\dim\mathfrak e_8
=248\\
&\Longrightarrow
\exists C,\quad C^{30}=I.
\end{aligned}
}
\]

\end{document}
